\section{Anchored antiderivatives and a finite derivative--primitive algebra}
\label{rlc:sec:primitives}

\subsection{The covariant fugacity ladder}
The transform immediately relates changing total order to integrating
or differentiating the logarithmic product argument.

\begin{theorem}[Covariant primitives and shift derivatives]
\label{rlc:thm:ladder}
For $\Rea a>0$, $W\in\C$, and real $\mu$,
\begin{align}
 (\partial_\mu+a-1)\CC_{r,a}(W;\mu)
      &=\CC_{r,a}(W-1;\mu),\label{rlc:eq:muladder}\\
 \CC_{r,a}(W+1;\mu)
      &=e^{-(a-1)\mu}\int_{-\infty}^{\mu}
              e^{(a-1)v}\CC_{r,a}(W;v)\dd v,
        \label{rlc:eq:muprimitive}\\
 \partial_a^n\CC_{r,a}(W;\mu)
      &=(-1)^n(W)_n\CC_{r,a}(W+n;\mu),\qquad n\ge0.
       \label{rlc:eq:aderivative}
\end{align}
The primitive in \eqref{rlc:eq:muprimitive} is fixed by the condition
$e^{(a-1)\mu}\CC_{r,a}(W+1;\mu)\to0$ as $\mu\to-\infty$.
All identities can be differentiated in $W$.
\end{theorem}
\begin{proof}
In \eqref{rlc:eq:barnes}, $\partial_\mu+a-1$ multiplies the integrand
by $p+a-1$, proving \eqref{rlc:eq:muladder}. Differentiating the power
$(p+a-1)^{-W}$ in $a$ gives \eqref{rlc:eq:aderivative}.
The contour bound for real negative $\mu$ is $O(e^{c\mu})$ on compact
order sets. Since $c+\Rea a-1>0$, multiplication by
$e^{(a-1)\mu}$ tends to zero at $-\infty$ and the integral in
\eqref{rlc:eq:muprimitive} converges. Integrating
\eqref{rlc:eq:muladder} at order $W+1$ proves the anchored formula.
The same bounds with logarithmic factors justify all order derivatives.
\end{proof}

At $a=1$, this is an ordinary primitive ladder in $\mu$, or an Euler
primitive ladder in $q=e^\mu$ with measure $\dd q/q$. Raising $W$ by
any positive integer supplies repeated primitives with all constants
fixed. At other shifts the integrating factor is indispensable.

Let $D_a=\partial_\mu+a-1$ and let $J_a$ denote the anchored integral
operator on the right of \eqref{rlc:eq:muprimitive}. On the linear span
of the convolution profiles and their order derivatives,
\begin{equation}\label{rlc:eq:commutators}
 [\partial_a,D_a]=1,\qquad
 [\partial_a,J_a]=-J_a^2.
\end{equation}
The first identity is immediate. For the second, differentiation of the
integrating factor gives
\[
 (\partial_aJ_a-J_a\partial_a)h(\mu)
   =-\int_{-\infty}^{\mu}(\mu-v)e^{-(a-1)(\mu-v)}h(v)\dd v
   =-J_a^2h(\mu).
\]
Absolute convergence justifies the repeated integral. Thus a derivative
in the shift does not commute with a covariant primitive; the correction
is an explicitly normalized second primitive.

Combining \eqref{rlc:eq:aderivative} with order differentiation gives
\begin{equation}\label{rlc:eq:jetderivative}
 \partial_a^n\partial_W^M\CC_{r,a}(W;\mu)
 =(-1)^n\sum_{h=0}^{\min(M,n)}\binom Mh
       \partial_W^h(W)_n\,
       \partial_W^{M-h}\CC_{r,a}(W+n;\mu).
\end{equation}
The coefficients at $W=0$ are precisely the finite harmonic
polynomials in \eqref{rlc:eq:Bellharmonic}. Equations
\eqref{rlc:eq:commutators}--\eqref{rlc:eq:jetderivative}, together with
the reciprocal harmonic product in \eqref{rlc:eq:completeharmonic},
provide a concrete compatible derivative--primitive algebra in this
family. This addresses a specified instance of the compatibility
question in the Mellin continuation \cite{MLD}, not an assertion that
every harmonic calculus in the repository has thereby been unified.

\subsection{Anchored integration in the Hurwitz shift}
For $\Rea a,\Rea b>0$, integration can be taken along any path in the
right half-plane. A pole-free fundamental identity is
\begin{equation}\label{rlc:eq:Eprimitive}
 \int_a^b\E_c(W)\dd c=\zeta(W,a)-\zeta(W,b).
\end{equation}
It follows from $\partial_c\zeta(W,c)=-W\zeta(W+1,c)$.
The right side is entire in $W$, since the residues at $W=1$ cancel.
The central even-depth analogue is
\begin{equation}\label{rlc:eq:evenprimitive}
 \int_a^b\CC_{2m,c}(W;0)\dd c
 =\frac1{(2m-1)!}\sum_{k=0}^{m-1}A_{m,k}(W)_{2k}
      [\zeta(W+2k,a)-\zeta(W+2k,b)].
\end{equation}
Every bracket is taken as an entire combined difference before
specialization.

At odd depth, put $\Delta\eta(v)=\eta(v,a)-\eta(v,b)$. Then
\begin{align}
 \int_a^b\CC_{2m+1,c}(W;0)\dd c
 =\frac1{(2m)!}\biggl[&B_{m,0}\frac{\Delta\eta(W-1)}{W-1}\notag\\
 &+\sum_{k=1}^{m}B_{m,k}(W)_{2k-1}
                          \Delta\eta(W+2k-1)\biggr].
 \label{rlc:eq:oddprimitive}
\end{align}
The first quotient is removable at $W=1$, because $\eta(0,a)=\eta(0,b)$.
Its value there is $\eta^{[1]}(0,a)-\eta^{[1]}(0,b)$.
These two formulas follow by differentiating their right sides and using
the anchor $b=a$; they therefore specify the integration constants as
well as the derivatives.

\subsection{Stieltjes antiderivatives and Gamma normalization}
Taking order derivatives of \eqref{rlc:eq:Eprimitive} at zero gives
\begin{equation}\label{rlc:eq:gammaprimitive}
 \int_a^b\gamma_0(c)\dd c=\log\Gamma(a)-\log\Gamma(b),
\end{equation}
and, for every $n\ge0$,
\begin{equation}\label{rlc:eq:stieltjesprimitive}
 \boxed{\displaystyle
 \int_a^b\gamma_n(c)\dd c
 =\frac{(-1)^n}{n+1}
       [\zeta^{[n+1]}(0,a)-\zeta^{[n+1]}(0,b)].}
\end{equation}
Equivalently an indefinite primitive is
$(-1)^{n+1}\zeta^{[n+1]}(0,a)/(n+1)$ plus a constant.
These are the classical Hurwitz order-derivative primitives, now attached
to the ordinary product identities of Theorem~\ref{rlc:thm:stieltjes}.
We do not count the familiar $n=0$ log-Gamma formula as a new result.

The odd-depth removable primitive also has a finite Gamma value:
\[
 \eta^{[1]}(0,a)=\log\Gamma(a/2)-\log\Gamma((a+1)/2)
                         -\tfrac12\log2.
\]
This follows from \eqref{rlc:eq:eta} and the classical value
$\zeta^{[1]}(0,a)$. In particular, the only potentially singular
coefficient in \eqref{rlc:eq:oddprimitive} is resolved explicitly by
Gamma differences rather than an unspecified integration constant.
